% SPDX-FileCopyrightText: 2026 Semih Burak, RWTH Aachen University
%
% SPDX-License-Identifier: Apache-2.0 WITH LLVM-exception
%
% From source to executable through the SPMD dialect. Build with `make` in
% this directory; the page is cropped to the picture, so the PDF can be
% \includegraphics'd as is, or the tikzpicture copied into a document that
% loads the same libraries and the styles below.

\documentclass{article}
\usepackage[T1]{fontenc}
\usepackage{lmodern}
\usepackage[scaled=0.82]{beramono}
\usepackage{tikz}
\usetikzlibrary{arrows.meta,backgrounds,calc,fit,positioning}
\usepackage[active,tightpage]{preview}
\PreviewEnvironment{tikzpicture}
\setlength\PreviewBorder{2pt}
% Wider than the picture, so its line is not reported as overfull.
\setlength\textwidth{50cm}

% Okabe-Ito: distinguishable with colour-vision deficiencies.
\definecolor{cFrontend}{HTML}{666666}
\definecolor{cRaise}{HTML}{0072B2}
\definecolor{cOpt}{HTML}{009E73}
\definecolor{cSelect}{HTML}{E69F00}
\definecolor{cLower}{HTML}{CC79A7}
\definecolor{cGeneric}{HTML}{56B4E9}
\definecolor{cCodegen}{HTML}{D55E00}

\tikzset{
  every picture/.style={font=\sffamily\footnotesize, >={Latex[length=1.8mm]}},
  % An IR level or other artifact the pipeline passes along.
  ir/.style={draw, thick, rounded corners=5pt, fill=white, align=center,
    inner sep=3pt, minimum height=8mm},
  % A pass or tool.
  pass/.style={draw=#1!80!black, fill=white, align=center, inner sep=2.5pt,
    font=\ttfamily\scriptsize},
  % A stage: the background box around its passes, and its title.
  stage/.style={draw=#1, fill=#1!12, rounded corners=3pt, inner sep=4pt},
  title/.style={font=\sffamily\scriptsize\bfseries, text=#1!65!black,
    inner sep=1pt},
  note/.style={font=\sffamily\scriptsize, align=center, inner sep=1pt},
  flow/.style={->, thick},
  sub/.style={->, draw=black!60},
  side/.style={->, densely dashed, draw=black!60},
  model/.style={pass=cLower, minimum width=15mm, font=\sffamily\scriptsize},
}

% Horizontal room for an arrow: from a stage box to the next node, and from a
% node to the first pass of the next stage (which also spans the box's padding).
\newlength\irgap \setlength\irgap{6mm}
\newlength\passgap \setlength\passgap{7.5mm}

% \stagebox{colour}{name}{fit list}: the box behind a stage's passes.
\newcommand\stagebox[3]{%
  \begin{scope}[on background layer]
    \node[stage=#1, fit=#3] (#2) {};
  \end{scope}}

\begin{document}
\begin{tikzpicture}

  % ---------------------------------------------------------------- row 1

  \node[ir] (src) at (0,0) {C/C++ program\\[1pt]
    {\scriptsize MPI $\cdot$ OpenSHMEM $\cdot$ NVSHMEM}\\
    {\scriptsize NCCL $\cdot$ CUDA $\cdot$ OpenMP}};

  % Frontend. A stage's widest pass is placed first, so that the box around
  % it keeps its distance from the node before.
  \node[pass=cFrontend, anchor=west, font=\sffamily\scriptsize]
    (polygeist) at ($(src.east) + (\passgap, -3mm)$)
    {Polygeist raising\\(Enzyme-JAX)};
  \node[pass=cFrontend, above=3mm of polygeist] (clang) {reactant-clang};
  \node[title=cFrontend, above=1.5mm of clang] (tFrontend) {Frontend};
  \draw[sub] (clang) -- (polygeist);
  \stagebox{cFrontend}{sFrontend}{(tFrontend) (clang) (polygeist)}

  \node[ir, anchor=west] (mlir) at ($(sFrontend.east |- src) + (\irgap, 0)$)
    {MLIR\\[1pt]{\scriptsize llvm, scf, affine,}\\{\scriptsize memref, omp}};

  % Raising
  \node[pass=cRaise, anchor=west] (cacts) at ($(mlir.east) + (\passgap, -3mm)$)
    {convert-apiCalls-to-spmd};
  \node[pass=cRaise, above=3mm of cacts] (omp2spmd) {omp-to-spmd};
  \node[title=cRaise, above=1.5mm of omp2spmd] (tRaise) {Raising};
  \draw[sub] (omp2spmd) -- (cacts);
  \stagebox{cRaise}{sRaise}{(tRaise) (omp2spmd) (cacts)}

  \node[ir, very thick, anchor=west, minimum width=17mm] (spmd)
    at ($(sRaise.east |- src) + (\irgap, 0)$) {SPMD IR};

  % What CACTS maps each programming model's calls with.
  \node[ir, thin, draw=black!60, rounded corners=1pt, anchor=south west,
    font=\sffamily\scriptsize] (apijson)
    at ($(sRaise.north west) + (0, 5mm)$) {API descriptions\\\texttt{api.json} (USIS)};
  \draw[side] (apijson) -- (apijson |- sRaise.north);

  % Optimization rewrites SPMD IR into SPMD IR, so it sits above it.
  \node[pass=cOpt, above=10mm of spmd] (coalesce) {coalesce-comm};
  \node[pass=cOpt, above=1.5mm of coalesce] (barriers)
    {eliminate-redundant-barriers};
  \node[pass=cOpt, above=1.5mm of barriers] (redcomm)
    {eliminate-redundant-comm};
  \node[pass=cOpt, above=1.5mm of redcomm] (deadcomm) {eliminate-dead-comm};
  \node[title=cOpt, above=1.5mm of deadcomm] (tOpt) {Optimization (optional)};
  \stagebox{cOpt}{sOpt}{(tOpt) (deadcomm) (redcomm) (barriers) (coalesce)}

  % ---------------------------------------------------------------- row 2

  % Model selection
  \node[note, anchor=north west] (models) at ($(src.west) + (1.5mm, -33mm)$)
    {user option:\\
     \texttt{models=shmem=mpi,}\\\texttt{\phantom{models=}omp=keep}};
  \node[pass=cSelect, above=1.5mm of models] (lowerspmd) {lower-spmd};
  \node[title=cSelect, above=1.5mm of lowerspmd] (tSelect) {Model selection};
  \stagebox{cSelect}{sSelect}{(tSelect) (lowerspmd) (models)}

  % Model-specific lowering
  \node[model, anchor=west] (lSHMEM)
    at ($(sSelect.east) + (\passgap, 0)$) {OpenSHMEM};
  \node[model, above=1.5mm of lSHMEM] (lMPI) {MPI};
  \node[model, below=1.5mm of lSHMEM] (lNVSHMEM) {NVSHMEM};
  \node[model, right=1.5mm of lSHMEM] (lCUDA) {CUDA};
  \node[model, above=1.5mm of lCUDA] (lNCCL) {NCCL};
  \node[model, below=1.5mm of lCUDA] (lOMP) {OpenMP};
  \node[note, font=\ttfamily\scriptsize, anchor=south west] (lowerto)
    at ($(lMPI.north west) + (0, 1.5mm)$)
    {lower-spmd-to-\textsf{\itshape model}};
  \node[title=cLower, above=1.5mm of lowerto.north west, anchor=south west]
    (tLower) {Model-specific lowering};
  \stagebox{cLower}{sLower}{(tLower) (lowerto) (lMPI) (lNVSHMEM) (lOMP)}

  \node[ir, anchor=west] (calls) at ($(sLower.east |- lSHMEM) + (\irgap, 0)$)
    {MLIR\\[1pt]{\scriptsize runtime calls,}\\{\scriptsize omp dialect}};

  % Generic lowering
  \node[pass=cGeneric, anchor=west] (polyllvm)
    at ($(calls.east) + (\passgap, -2mm)$) {convert-polygeist-\\to-llvm};
  \node[title=cGeneric, above=1.5mm of polyllvm] (tGeneric)
    {Generic lowering};
  \stagebox{cGeneric}{sGeneric}{(tGeneric) (polyllvm)}

  \node[ir, anchor=west] (llvmdialect)
    at ($(sGeneric.east |- calls) + (\irgap, 0)$) {LLVM\\dialect};

  % Code generation
  \node[pass=cCodegen, anchor=west] (translate)
    at ($(llvmdialect.east) + (\passgap, 6mm)$) {mlir-translate};
  \node[pass=cCodegen, below=4mm of translate] (compile) {clang -c};
  \node[pass=cCodegen, below=4mm of compile] (link) {link};
  \node[title=cCodegen, above=1.5mm of translate] (tCodegen)
    {Code generation};
  \draw[sub] (translate) -- node[note, right=0.5mm] (nIR) {LLVM IR} (compile);
  \draw[sub] (compile) -- node[note, right=0.5mm] (nObj) {object} (link);
  \stagebox{cCodegen}{sCodegen}{(tCodegen) (translate) (link) (nIR) (nObj)}

  \node[ir, very thick, anchor=west] (exe)
    at ($(sCodegen.east |- calls) + (\irgap, 0)$) {Executable};

  % ------------------------------------------------------------------ flow

  \draw[flow] (src) -- (src -| sFrontend.west);
  \draw[flow] (mlir -| sFrontend.east) -- (mlir);
  \draw[flow] (mlir) -- (mlir -| sRaise.west);
  \draw[flow] (spmd -| sRaise.east) -- (spmd);

  % The optimization loop: out of SPMD IR, through the passes, back in.
  \draw[flow] ([xshift=-3mm]spmd.north) -- ([xshift=-3mm]spmd.north |- sOpt.south);
  \draw[flow] ([xshift=3mm]spmd.north |- sOpt.south) -- ([xshift=3mm]spmd.north);

  % Wrap from the end of row 1 to the start of row 2.
  \coordinate (wrap) at ($(spmd.south |- sFrontend.south)!0.5!(spmd.south |- sLower.north)$);
  \draw[flow, rounded corners=4pt] (spmd.south) -- (wrap)
    -- (wrap -| sSelect.north) -- (sSelect.north);

  \draw[flow] (sSelect.east |- lSHMEM) -- (sLower.west |- lSHMEM);
  \draw[flow] (sLower.east |- calls) -- (calls);
  \draw[flow] (calls) -- (calls -| sGeneric.west);
  \draw[flow] (llvmdialect -| sGeneric.east) -- (llvmdialect);
  \draw[flow] (llvmdialect) -- (llvmdialect -| sCodegen.west);
  \draw[flow] (exe -| sCodegen.east) -- (exe);

\end{tikzpicture}
\end{document}
